\chapter{Conclusion}
\label{chap:conclusion}

This chapter summarizes the contributions of this thesis, discusses limitations and challenges, and proposes directions for future work.

\section{Summary of Contributions}

This thesis presented an automated grading system for Entity-Relationship Diagrams that successfully addresses the key challenges of manual assessment in database design education. The main contributions are:

\textbf{Technical Contributions:}
\begin{itemize}
\item Automated ERD parsing system from Draw.io SVG files
\item Novel JSON prompting methodology for consistent AI grading
\item Multi-layer architecture with modular component separation
\item Comprehensive error handling and retry mechanisms
\end{itemize}

\textbf{Educational Contributions:}
\begin{itemize}
\item Reduced grading time from 15-20 minutes to under 10 seconds (110-146x speedup)
\item Detailed category-level feedback with actionable suggestions
\item Consistent, deterministic grading eliminating human fatigue effects
\item Scalability for large classes without compromising quality
\end{itemize}

\textbf{Research Contributions:}
\begin{itemize}
\item Demonstrated AI effectiveness for diagram assessment (100\% within ±5 points accuracy)
\item Validated reference calibration approach for AI scoring consistency
\item Established metrics for automated ERD grading quality
\end{itemize}

\section{Limitations and Challenges}

Current limitations include dependency on a single model answer, Draw.io format requirements, internet/API availability, and limited contextual understanding. Implementation challenges involved handling ERD notation variations, diverse cardinality formats, AI response consistency, and relationship participant extraction.

\section{Future Work}

Short-term enhancements include multi-format support, interactive feedback, partial credit mechanisms, and batch processing. Long-term directions involve fine-tuned AI models, explainable AI integration, adaptive learning systems, and real-time collaborative grading during diagram creation.

\section{Final Remarks}

This research demonstrates that large language models can be effectively applied to visual diagram assessment through appropriate parsing and structured prompting. The system provides a valuable tool for database design education while establishing a foundation for automated assessment of other visual artifacts in computer science education.
